\section{Resultados por indicador}
\label{sec:resultados}

A Figura~\ref{fig:contrasteoss} resume o contraste central da pesquisa:
quanto cada indicador difere entre a gestão por OSS e a administração
Direta, já descontados complexidade, porte, perfil de longa
permanência e ano, com as margens de erro agrupadas por hospital.

O quadro que emerge é o de um \textbf{regime operacional claramente
mais intenso nas OSS}: produção de saídas \textbf{45,4\% maior},
ocupação de internação \textbf{23,6\% maior} e uso da UTI, entre os
hospitais que a possuem, \textbf{20,4\% maior}, todos com significância
estatística confortável. É uma diferença de regime entre redes com
missões distintas, não um efeito demonstrado da conversão de um
hospital individual (ver as Seções~\ref{sec:efeitogestao}
e~\ref{sec:inferenciafinal}). O \textbf{tempo de permanência é 10,2\%
menor} nas OSS na comparação de cada hospital consigo mesmo (p igual a
0,021, valor que, agrupado por hospital com apenas cinco tratados,
tende a ser otimista; a Seção~\ref{sec:inferenciafinal} o submete a
permutação, que devolve p igual a 0,153), o que fecha a conta do giro:
mais pacientes passam pelos mesmos
leitos porque cada um fica menos tempo internado. O \textbf{custo real
por saída} aparece 6,7\% menor, mas com margem de erro grande demais
para afirmar diferença (o intervalo vai de 22,6\% de queda a 12,4\% de
alta): com apenas cinco hospitais mudando de gestão, os dados não
permitem cravar se as OSS são mais baratas ou mais caras por paciente
atendido. Os efeitos de ano do mesmo modelo confirmam, agora em
ambiente controlado, a queda real de custo da década: 2025 fica 11,6\%
abaixo de 2015 em preços comparáveis.

Na \textbf{mortalidade}, o modelo Beta com inflação em zero e
especificação de Mundlak estima, para a mudança de Direta para OSS
dentro do mesmo hospital, uma redução equivalente a \textbf{1,12 ponto
percentual} na taxa (p igual a 0,085, significância marginal); esse
número não sobrevive intacto ao confronto com o controle sintético da
Seção~\ref{sec:inferenciafinal} e não deve ser citado isoladamente. Duas
leituras complementares importam. Primeiro, o termo de Mundlak (que
captura a diferença estrutural entre hospitais OSS e os demais) é
positivo e não significativo: as OSS não são hospitais de mortalidade
estruturalmente baixa; a queda estimada vem da trajetória dos
conversores. Segundo, o componente de zeros do modelo confirmou um
fato notável da primeira parte: \textbf{nenhum hospital OSS ou
Municipal registra ano sem óbito} no painel, de modo que a chance de
mortalidade zero foi modelada apenas para Filantrópicos e Privados
contra o bloco público. Na \textbf{fração de alta complexidade}, o
modelo ZOIB não encontra diferença entre OSS e Direta em nenhum dos
componentes; quem explica essa fração é a complexidade estrutural (nos
zeros) e o porte (na média), como esperado.

\begin{figure}[H]
\centering
\includegraphics[width=.96\textwidth]{fig_est_03_contraste_oss_direta.png}
\caption{Contraste OSS vs Direta nos modelos principais, em efeito
percentual com intervalo de 95\%; a mortalidade, que tem escala
própria, aparece anotada no título. Custo e TMP: identificação
\emph{within} (efeitos fixos, sustentada pelos cinco conversores);
produção e ocupações: contraste predominantemente \emph{between} entre
as redes OSS e Direta, de natureza associativa.}
\label{fig:contrasteoss}
\end{figure}

\comoler{cada linha é um indicador; o quadrado verde é a melhor
estimativa do efeito de ser OSS em vez de Direta, e a haste horizontal
é a margem de erro. Quando a haste cruza a linha vertical do zero, como
no custo, não se pode afirmar diferença; quando fica inteira de um
lado, como na produção e na ocupação, a diferença é estatisticamente
firme. Valores à direita do zero indicam OSS maior; à esquerda, OSS
menor.}
